% =============================================================
% Chapitre 6 --- Validation, résultats et démonstration
% Source : ../../CH6_Validation_mesures_et_evaluation.md
% Ancrage des requalifications : ../../DOSSIER_TECHNIQUE_MENAL.md
% Annexes associées : appendices/annexe_e_protocoles_tests.tex,
%                     appendices/annexe_f_evaluation_semantique.tex
% Huit sections (plafond respecté) :
%   1. Stratégie de validation
%   2. Résultats de la campagne
%   3. Détection et enrichissement
%   4. Performance, résilience et coût
%   5. Démonstration de bout en bout et conformité aux référentiels
%   6. Ce que le filtrage applicatif ne protège pas
%   7. Du socle au service
%   8. Conditions de validité et réponse aux objectifs
% =============================================================
\chapter{Validation, résultats et démonstration}
\label{ch:validation}
\soustitrechapitre{Campagne de tests adverses, évaluation de la détection et mesures}
% Resserrement local du placement des flottants (chapitre dense en tableaux).
%     Valeurs restaurees en fin de chapitre pour ne pas affecter la suite du document.
\setcounter{topnumber}{3}\setcounter{bottomnumber}{2}\setcounter{totalnumber}{5}
\renewcommand{\topfraction}{0.92}\renewcommand{\bottomfraction}{0.7}
\renewcommand{\textfraction}{0.06}\renewcommand{\floatpagefraction}{0.78}
\setlength{\LTpre}{6pt}\setlength{\LTpost}{6pt}

Les chapitres précédents ont analysé, conçu et construit ; celui-ci mesure : une architecture Zero
Trust ne se déclare pas, elle se démontre \cite{nist800207}. Les vingt tests de la
\cref{sec:campagne} sont ceux dérivés de l'analyse de risque du \cref{ch:besoins-menaces}, et chaque mesure
rapportée ici est soit une valeur relevée et datée, soit une propriété établie par une preuve
nommée, jamais une estimation présentée comme mesurée. Le chapitre suit quatre axes --- sécurité,
détection, performance et résilience --- puis énonce la frontière du service rendu.

\section{Stratégie de validation}
\label{sec:strategie}

Chaque protocole reçoit l'un de quatre résultats, appliqués sans exception, y compris
défavorables. \textbf{Conforme} : le critère d'acceptation est couvert \emph{intégralement} par
une preuve nommée --- le plus souvent un refus effectivement observé, en direct ou par un test
automatisé rejoué à chaque livraison ; par exception, un inventaire exhaustif lorsque le critère
est lui-même un inventaire. \textbf{Partiellement conforme} : le critère est couvert en partie par
une preuve nommée, la part non couverte étant écrite comme une borne du dispositif de mesure.
\textbf{Non conforme} : le critère n'est pas tenu ; l'exigence associée redevient ouverte et le
scénario de menace du \cref{ch:besoins-menaces} se rouvre avec elle. \textbf{Planifié} : le protocole relève du
plan de validation continue (\cref{sec:plan-validation}). \textbf{Aucun protocole ne porte ce
résultat} à la date de \dateref{} ; la convention est néanmoins conservée telle qu'elle a été
publiée --- la retirer après coup reviendrait à ajuster la méthode au résultat obtenu.

Cette convention repose sur une hiérarchie de preuves déclarée à l'avance, qui tranche les cas
limites. Rang 1, le \textbf{test automatisé rejoué à chaque livraison} : il établit la propriété
\emph{et} sa persistance, un tiers peut le contredire en relançant la chaîne, et il ne vaut que
pour le chemin qu'il emprunte. Rang 2, le \textbf{refus réellement survenu en exploitation}, plus
fort sur le fond puisque le contrôle a arrêté un cas non provoqué, mais non répétable à volonté.
Rang 3, la \textbf{vérification en direct, datée}, qui établit un état à un instant et ne dit rien
de la dérive ultérieure. Rang 4, la \textbf{configuration déclarée en code et vérifiée sur la
plateforme}, qui établit ce que le système \emph{est}, non ce qu'il \emph{fait} sous contrainte. La
capture d'écran ne figure pas dans cette échelle : elle restitue une preuve d'un des quatre rangs,
elle n'en est jamais une. Les sept cadres de capture du mémoire ---
\cref{fig:capture-K35,fig:capture-K36,fig:capture-K23,fig:capture-K14,fig:capture-K05,fig:capture-K33,fig:capture-K34}
--- restituent visuellement des preuves de ces rangs, nommées et datées dans leur légende ;
conformément à la hiérarchie ci-dessus, \textbf{ils illustrent des preuves nommées et n'en
tiennent jamais lieu}. Chacun porte, tant que l'image n'est pas déposée, la spécification de la
prise : la vue attendue, le contenu qui vaut preuve et ce qui doit être masqué.

\begin{alertbox}
Une preuve de rang 4 ne suffit jamais à prononcer « conforme », et le projet en a fait lui-même
la démonstration : la porte d'analyse statique de la chaîne de livraison applicative était
déclarée bloquante, correctement écrite, et n'évaluait plus aucune règle depuis au moins deux
semaines sans jamais faire échouer la livraison (\cref{subsec:incidents}). Une porte déclarée n'est pas une porte
vérifiée --- d'où plusieurs protocoles dont le dispositif est en place et qui restent ici
\emph{partiellement} conformes.
\end{alertbox}

Les mesures sont réalisées sur l'environnement de recette, représentatif de la cible décrite mais
sans trafic d'utilisateurs réels (\cref{sec:cloture}). Les preuves ponctuelles sont datées jusqu'au
25/08/2026 ; les preuves automatisées sont rejouées à chaque livraison, sans date unique. Les
protocoles, leurs critères d'acceptation formulés avant exécution et leurs réserves sont consignés
en annexe~\ref{ann:protocoles-tests}.

\section{Résultats de la campagne de tests}
\label{sec:campagne}

Les vingt tests dérivés du \cref{ch:besoins-menaces} couvrent les quatre axes, le protocole T\emph{n} validant
l'exigence EX\emph{n} de même rang (\cref{tab:matrice-tracabilite}). Chaque ligne du
\cref{tab:synthese-campagne} nomme la preuve invoquée et sa date, ou son mode d'exécution
automatisée ; le protocole complet, son critère d'acceptation et la portée exacte de la preuve se
lisent test par test en annexe~\ref{ann:protocoles-tests}.

{\footnotesize
\setlength{\tabcolsep}{3pt}
\begin{longtable}{|P{0.8cm}|P{1.55cm}|P{12.9cm}|}
\caption[Synthèse de la campagne de vingt tests]{Synthèse de la campagne de vingt tests, à la date de \dateref{}. Le protocole
T\emph{n} valide l'exigence EX\emph{n} de même rang ; « Partiel » : partiellement conforme}
\label{tab:synthese-campagne} \\
\hline
\textbf{Test} & \textbf{Résultat} & \textbf{Preuve invoquée} \\
\hline
\endfirsthead
\hline
\textbf{Test} & \textbf{Résultat} & \textbf{Preuve invoquée} \\
\hline
\endhead
T1 & \textbf{Non conforme} & Injections SQL, script inter-sites et inclusion de fichier
refusés en~403 le 23/08/2026 ; motif brut de traversée de chemin en~302, alors que le critère
exige le refus de \emph{toutes} les requêtes (\cref{sec:non-conformite-t1}) \\ \hline
T2 & Conforme & Refus observé le 19/08/2026 : dix réponses~422 puis une~429 à la onzième
tentative \\ \hline
T3 & Partiel & Balayage tracé, journaux du répartiteur collectés sans filtre.
\emph{Réserve} : essai de charge non documenté \\ \hline
T4 & Partiel & Inventaire : aucune adresse publique, accès par voie privée, transport
chiffré imposé. \emph{Réserve} : échec de connexion externe déduit de l'absence de route \\ \hline
T5 & Conforme & Deux rôles exacts et usurpation refusée le 23/08/2026 ; test automatisé de
lecture interdite d'un secret rejoué à chaque livraison \\ \hline
T6 & Partiel & Test automatisé d'appel sans jeton, en porte bloquante ; entrée interne
exclusive et invocation nominative. \emph{Réserve} : l'appel vient de l'extérieur du réseau
privé \\ \hline
T7 & Conforme & Le critère est un inventaire, et il est vide : aucune clé de compte de
service, fédération contrainte sur le dépôt \emph{et} la branche. Garantie vérifiée, non
préventive (écart É7) \\ \hline
T8 & \textbf{Conforme} & \textbf{Refus réellement survenu} : la porte d'analyse a bloqué une
livraison sur une vulnérabilité critique de \texttt{node-tar} 6.2.1, embarquée par l'image de base
(\cref{sec:preuve-t11}) \\ \hline
T9 & Partiel & Registre unique, résolution par empreinte sur le chemin de description en
code. \emph{Réserve} : aucun refus opposé ; application hébergée déployée par étiquette flottante
(écart É8) \\ \hline
T10 & Partiel & Porte de détection de secrets bloquante à chaque livraison, sur
l'historique complet ; indépendance des secrets établie par inventaire. \emph{Réserve} : aucun
refus observé \\ \hline
T11 & \textbf{Conforme} & \textbf{Test de bout en bout automatisé} tentant réellement
l'écriture interdite dans la table de preuves sous l'identité d'enrichissement, rejoué à chaque
livraison ; droits vérifiés table par table le 19/08/2026 (\cref{sec:preuve-t11}) \\ \hline
T12 & Partiel & Alerte d'absence de trafic reçue le 22/08/2026, 36~min~19~s après
l'interruption. \emph{Réserve} : une route de journaux coupée 1~h~08 le même jour n'a ouvert aucune
alerte et 1\,148 requêtes n'ont laissé aucune trace \\ \hline
T13 & Partiel & Journalisation d'audit des accès aux données activée en lecture et en
écriture sur quatre services ; activité d'administration non désactivable. \emph{Réserve} :
rétention de trente jours non verrouillée, sans acheminement dédié \\ \hline
T14 & \textbf{Non conforme} & Le 25/08/2026, une sonde portant la configuration réseau de la
révision déployée atteint une destination externe en 412~ms ; le refus par défaut en sortie,
vérifié le même jour sur le sous-réseau, ne s'applique pas au chemin du composant
(\cref{sec:non-conformite-t1}) \\ \hline
T15 & Partiel & 22/08/2026 : 350 détections traitées en 7 exécutions, entrées égales aux
sorties, durée maximale 96~s, coût dominé par le forfait de balayage. \emph{Réserve} : 150 des 500
détections injectées sortent de la fenêtre de deux heures avant d'être lues, sans erreur ni
trace \\ \hline
T16 & Partiel & Mécanisme sollicité en conditions réelles le 08/08/2026, planification sans
différence le 19/08/2026, contrôle quotidien de dérive de l'isolation. \emph{Réserve} : aucune
dérive planifiée sur l'infrastructure décrite en code \\ \hline
T17 & Conforme & Quatre refus opposés le 20/08/2026 à trois identités de service, chacun
nommant le droit manquant, avec contrôle positif. \emph{Réserve} : emplacement de l'état
provisionné hors description en code \\ \hline
T18 & Partiel & Revue de fusion et application manuelle par l'administrateur.
\emph{Réserve} : aucune approbation outillée, aucun environnement protégé déclaré \\ \hline
T19 & Partiel & Deux notifications de budget reçues le 21/08/2026, seuils 90\,\% et
100\,\%, sur l'alerte déclarée dans le code d'infrastructure le 20/08/2026. \emph{Réserve} : délai
compris entre 5~h~48 et 8~h~18, et la notification ne nomme pas le poste de consommation \\ \hline
T20 & \textbf{Conforme} & Le 24/08/2026, trois sorties refusées --- externe, croisée entre
locataires, et vers la base --- avec les trois lignes de refus retrouvées dans l'entrepôt en
2~min~55~s à 4~min~12~s, les horodatages du refus et de l'échec client se recoupant à 10,3~s
près \\
\hline
\end{longtable}
}

\textbf{Bilan à la date de \dateref} : \textbf{7 conformes, 11 partiellement conformes et 2 non
conformes}. Les vingt protocoles portent une preuve nommée, vérifiable et rattachée à une date ou à
une exécution automatisée ; le plan de validation continue (\cref{sec:plan-validation}) ne porte
plus des protocoles sans résultat, mais la cadence à laquelle six d'entre eux sont rejoués.

Ce bilan appelle une lecture méthodologique : les onze résultats partiels ne sont pas des
demi-échecs, ce sont des propriétés établies dont une part précise du critère --- le plus souvent
la tentative négative depuis l'extérieur du périmètre, borne du dispositif de mesure énoncée en
\cref{sec:limites} --- n'a pas encore été mise à l'épreuve, et cette part est nommée ligne
à ligne.

\begin{keybox}[title={Une campagne sans non-conformité serait suspecte}]
Une campagne de vingt tests dont aucun n'échouerait indiquerait que les tests ont été choisis
après coup, pour tomber juste. Ceux-ci ont été dérivés de l'analyse de risque du \cref{ch:besoins-menaces},
\emph{avant} la conception de l'architecture, et publiés avec leur critère d'acceptation.
Les deux non-conformités, T1 et T14, sont la conséquence directe de ce choix de méthode : elles
sont le prix, et la preuve, de l'antériorité des critères. T14 en est l'illustration la plus
nette --- son critère, écrit avant la refonte réseau, aurait été trivialement satisfait par une
reformulation après coup en « absence de chemin public \emph{vers} le composant ».
\end{keybox}

\subsection{Les deux non-conformités assumées : T1 et T14}
\label{sec:non-conformite-t1}

Le protocole T1 exige que \emph{toutes} les charges d'attaque correspondant aux familles de
règles activées au périmètre soient rejetées avant d'atteindre la charge de travail. Le
23/08/2026, les charges d'injection SQL, de script inter-sites et d'inclusion de fichier local
ont bien reçu un refus en~403 ; une forme brute de traversée de chemin a reçu une redirection
302. Le critère n'est pas tenu, et le résultat est prononcé non conforme.

Le mécanisme n'est pas élucidé, et le mémoire le dit plutôt que de lui inventer une cause.
L'hypothèse de travail est une normalisation d'URL effectuée en amont de l'évaluation des règles :
la requête serait réécrite, puis redirigée, avant que le moteur de filtrage n'ait à se prononcer
sur la forme brute qui devait le déclencher. Cette hypothèse est vérifiable, et sa vérification
fait partie du correctif.

Le correctif est identifié et son ordre d'application est arrêté : activation des jeux de règles
d'anomalie de protocole et d'application de méthode, d'abord en observation seule pour mesurer le
volume de faux positifs sur le trafic de recette, puis en blocage une fois ce volume connu.
L'ordre importe : passer directement en blocage sur un jeu générique échangerait une lacune de
couverture contre une indisponibilité. Tant que ce correctif n'est pas appliqué et T1 rejoué,
l'exigence EX1 reste ouverte et le scénario de menace SO1 avec elle.

\textbf{T14 --- une non-conformité produite par la refonte, non masquée par elle.} Le protocole
exige l'échec d'une connexion sortante émise depuis le composant d'encodage. Le 25/08/2026, deux
sondes ne différant que par leur chemin de sortie ont été exécutées sous l'identité du composant :
celle qui emprunte le sous-réseau et l'étiquette du composant a été refusée deux fois par la règle
de refus par défaut en sortie, avec ses deux lignes de journal ; celle qui emprunte le chemin de la
révision réellement déployée --- sortie managée par la plateforme d'exécution, sans accès au réseau
privé --- a reçu un code~204 en 412~ms. La règle refuse ; le composant n'est pas derrière elle.
Le composant a été retiré du sous-réseau les 12 au 14/08 au motif, exact, qu'il n'a aucune
dépendance privée : ce motif justifiait le retrait pour la sortie \emph{nécessaire}, non pour la
sortie \emph{possible}. Le correctif est identifié et son coût est d'une révision sans changement
applicatif --- réintégrer le composant au sous-réseau et lui poser une étiquette pour laquelle
aucune règle d'autorisation de sortie n'existe --- au prix de le rendre à nouveau tributaire de la
disponibilité du sous-réseau. Il vaut identiquement pour l'interface web de l'application hébergée,
seconde charge de travail à sortie managée non gouvernée. Tant qu'il n'est pas appliqué et T14
rejoué, l'exigence EX14 reste ouverte et le scénario de menace SO17 avec elle.

\subsection{Deux preuves déterminantes et deux réserves levées}
\label{sec:preuve-t11}

Deux des sept résultats conformes reposent sur des preuves d'une nature que le mot « test »
recouvre mal, et qui portent plus loin qu'une capture d'écran.

\textbf{T11 --- le meilleur test de la campagne.} Le principe cardinal de la chaîne de détection
est que le moteur d'enrichissement ne doit pas pouvoir modifier les preuves qu'il analyse. Un
test de bout en bout automatisé, rejoué à chaque livraison, le vérifie : sous l'identité de la
tâche d'enrichissement obtenue par usurpation contrôlée, il exécute d'abord une lecture des
détections --- qui doit réussir, puisque la tâche lit pour enrichir --- puis tente réellement
l'insertion d'une ligne dans la table de preuves, et n'est vert que si cette insertion est
refusée.

Aux deux propriétés du rang~1 il en ajoute une troisième, qu'aucune image ne possède : il a une
histoire. Il échouait tant que la tâche d'enrichissement partageait son identité avec les requêtes
planifiées de détection, qui doivent, elles, écrire dans cette table ; \textbf{la séparation des
identités a été faite \emph{parce que} le test échouait}. Un test qui a déjà fait corriger
l'architecture qu'il vérifie mesure quelque chose. Sa portée s'énonce aussi précisément : il couvre
la table des détections, en lecture autorisée et en écriture refusée, le refus sur les journaux
bruts et le cas de la suppression étant établis par la vérification des droits table par table du
19/08/2026 ; et il s'abstient lorsque l'identité courante n'a pas le droit d'usurper l'identité
restreinte --- choix préféré à un vert obtenu avec les droits de l'opérateur, qui aurait mesuré la
mauvaise identité.

\textbf{T8 --- un refus qui n'a pas été provoqué.} La porte d'analyse de vulnérabilité de la
chaîne de livraison applicative a réellement arrêté une livraison, sur une vulnérabilité critique
de la bibliothèque d'archivage \texttt{node-tar} en version 6.2.1. Cette bibliothèque n'était pas
une dépendance du code écrit : elle était embarquée par le gestionnaire de paquets présent dans
l'image de base du tableau de bord, et le correctif retenu est instructif --- retirer ce
gestionnaire de l'image finale plutôt que tenter de mettre à jour une dépendance transitive dont
le projet n'est pas propriétaire. Ce refus n'a pas été mis en scène, il est survenu sur le code
de l'auteur, et il a modifié la façon de construire les images. La réserve qui l'accompagne ---
l'image est publiée au registre \emph{avant} d'être analysée, de sorte que la porte bloque le
déploiement et non la publication --- est établie avec son correctif en
\cref{sec:chaine-livraison}.

\textbf{Deux réserves de conception levées en cours de campagne.} La tâche d'enrichissement a
reçu une identité dédiée, limitée en lecture au jeu de données de supervision et en écriture à sa
seule table de résultat --- correction vérifiée par T11 et par la vérification des droits de T5 ;
T14 relève d'une autre propriété, que cette correction ne touche pas
(\cref{sec:non-conformite-t1}). La journalisation des refus du pare-feu a été activée et vérifiée
le 19/08/2026, ce qui clôt l'écart É1 : c'était la condition de recevabilité de T20, relevé le
24/08/2026 une fois le refus par défaut en sortie posé.

\subsection{Plan de validation continue}
\label{sec:plan-validation}

Six protocoles --- T12, T14, T15, T17, T19 et T20 --- portent un premier relevé daté entre le 20 et
le 25/08/2026 et une \textbf{cadence de rejeu}. Le plan est distinct de la campagne par sa nature :
celle-ci établit qu'un dispositif tient au moment où il est construit, celui-là qu'il tient encore
après avoir évolué. Ces six-là y figurent parce que la propriété qu'ils vérifient est celle qui se
dégrade --- un seuil qui dépend du volume, une règle de pare-feu qu'une révision peut contourner, un
droit sur un emplacement provisionné hors code : le plan est structuré par la fréquence de dérive du
contrôle vérifié, non par la commodité (\cref{tab:plan-validation}).

{\footnotesize
\setlength{\tabcolsep}{3pt}
\begin{longtable}{|P{0.8cm}|P{3.9cm}|P{1.6cm}|P{1.85cm}|P{6.7cm}|}
\caption[Plan de validation continue]{Plan de validation continue : premier relevé et cadence de rejeu}
\label{tab:plan-validation} \\
\hline
\textbf{Test} & \textbf{Propriété vérifiée} & \textbf{Premier relevé} & \textbf{Cadence de rejeu} &
\textbf{Déclencheur additionnel} \\
\hline
\endfirsthead
\hline
\textbf{Test} & \textbf{Propriété vérifiée} & \textbf{Premier relevé} & \textbf{Cadence de rejeu} &
\textbf{Déclencheur additionnel} \\
\hline
\endhead
T12 & Détection d'une interruption de collecte & 22/08/2026 & Trimestrielle & Toute modification
d'un acheminement de journaux \\ \hline
T14 & Absence de sortie réseau du composant d'inférence & 25/08/2026 & Trimestrielle & Application
du correctif : réintégration du composant au sous-réseau \\ \hline
T15 & Tenue de la chaîne d'enrichissement sous charge & 22/08/2026 & Trimestrielle & Changement de
modèle, de taille de lot ou de cadence de la tâche \\ \hline
T17 & Protection de l'état de l'infrastructure & 20/08/2026 & Semestrielle & Tout changement
d'emplacement ou de droits sur l'état \\ \hline
T19 & Détection d'une consommation anormale & 21/08/2026 & Trimestrielle & Toute révision du
montant ou des seuils du budget déclaré \\ \hline
T20 & Restriction des sorties réseau & 24/08/2026 & Trimestrielle & Ajout d'un locataire ou d'une
destination externe autorisée \\
\hline
\end{longtable}
}

Le premier rejeu de l'ensemble est fixé au quatrième trimestre 2026. Les trois prérequis identifiés
avant la campagne ont été levés, et leurs dates comptent autant que les résultats : refus de sortie
par défaut posé le 23/08/2026, alerte de budget déclarée en code le 20/08/2026, journalisation des
refus réseau acquise depuis le 19/08/2026. L'attendu des six avait été écrit avant exécution
(annexe~\ref{ann:protocoles-tests}) et confronté au relevé sans être ajusté : c'est ce qui donne
son poids au seul de ces six qui échoue.

\section{Évaluation de la chaîne de détection et de l'enrichissement}
\label{sec:eval-detection}

\subsection{Couverture obtenue}

Les deux précautions qu'appelle toute mesure de couverture --- dénominateur restreint aux
plateformes concernées, taux rapporté à une version précise du référentiel --- sont établies en
\cref{sec:attck-matrice} \cite{mitre_attack,mitre_attack_design2018}. Ce que la mesure ajoute ici
est l'\textbf{empreinte de l'instantané employé} : 872~vecteurs couvrant 697~techniques et
sous-techniques réparties sur \textbf{quinze tactiques}, dont TA0005 renommée « Stealth » et TA0112
« Defense Impairment », issue du découpage de l'ancienne « Defense Evasion » --- de quoi retrouver
la version Entreprise employée et reproduire le dénominateur. \emph{Réserve} : le script de
chargement tire le lot depuis la branche courante du dépôt public de MITRE, sans épinglage, de
sorte que l'instantané est daté par son chargement et non fixé par une référence immuable ; le
tableau de bord joint par \emph{nom} de tactique, ce qui rend un renommage visible plutôt que
muet.

Sept règles de détection sont actives (R1 à R7, chacune adossée à une requête planifiée), plus
cinq requêtes de normalisation. Faute d'inventaire versionné des techniques applicables, ni taux
ni graphique de couverture ne sont publiés : le nombre de règles actives n'est pas assimilé à un
taux de couverture ATT\&CK.

\subsection{Ce que le système ne détecte pas}

L'objectif O4 engage explicitement le mémoire à publier ce que le socle ne détecte pas. Le
\cref{tab:lacunes} tient cet engagement, cause par cause. Deux de ces lacunes sont
structurelles et se retrouvent en \cref{sec:limites-filtrage} : la chaîne de détection
n'observe que des signaux de transport, et une seule personne traduit les règles.

\begin{table}[!htbp]
\caption[Causes connues de non-couverture]{Causes connues de non-couverture, à la date de \dateref}
\label{tab:lacunes}
\centering
\footnotesize
\setlength{\tabcolsep}{3pt}
\begin{tabular}{|P{3.1cm}|P{4.5cm}|P{7.7cm}|}
\hline
\textbf{Cause} & \textbf{Origine} & \textbf{Conséquence} \\
\hline
Refus réseau tracés mais non qualifiés & Refus en sortie et refus croisés journalisés et
retrouvables (T20, 24/08/2026), qu'aucune des sept règles ne lit & Un mouvement latéral refusé est
\emph{tracé}, il n'est pas \emph{détecté} : ni détection, ni score, ni ligne au tableau de
bord \\ \hline
Plafond de traduction manuelle & Une seule personne traduit les règles publiques & Couverture
bornée par le volume maintenable ; exhaustivité non revendiquée \\ \hline
Absence de trafic réel & Environnement de recette sans usage réel & Seuils comportementaux non
calibrables \\ \hline
Exfiltration & Aucune règle dédiée au canal sortant ; seuls les refus sont journalisés, par
décision de coût, et deux charges conservent une sortie managée non gouvernée (T14, 25/08/2026) &
Un transfert sortant \emph{réussi} vers une destination autorisée n'est qualifié par rien, et
l'autorisation qui rend l'entrepôt joignable rend joignable tout entrepôt du même
fournisseur \\ \hline
Abus des droits et du plan de contrôle & Règles concentrées sur le filtrage applicatif et le
transport & Élévation de privilège et modification de politique d'accès non couvertes \\ \hline
Logique métier de l'application hébergée & La chaîne n'observe pas les actions fonctionnelles de
l'application & Un abus commis par un compte légitime reste invisible
(\cref{sec:limites-filtrage}) \\ \hline
Seuils codés en dur & Paramètres intégrés aux requêtes & Adaptation au volume réel lente et peu
traçable \\ \hline
Absence de test unitaire par règle & Validation surtout par scénarios de bout en bout &
Régression d'une règle difficile à isoler \\
\hline
\end{tabular}
\end{table}

\subsection{Qualité du signal, délai, et enrichissement sémantique}
\label{sec:eval-enrichissement}

Le volume d'alertes, le taux de faux positifs et les délais médian et maximal ne sont pas établis
par une campagne suffisamment large ; le scénario du 19/08 fournit une observation ponctuelle. La
règle R2, exécutée toutes les cinq minutes sur une fenêtre de quinze minutes, a produit une alerte
environ quinze minutes après les requêtes, et la tâche d'enrichissement s'exécute toutes les
15~min --- cadences de planification, non bornes garanties du délai de bout en bout. Le délai tient
à la nature de la règle (\cref{sec:chaine-detection}) : R2 agrège avant de décider et sa fenêtre
absorbe la latence d'ingestion, si bien que le délai attendu est voisin de la fenêtre et non de la
cadence de réévaluation. Une conséquence borne l'interprétation de tout délai mesuré : aucune
politique d'alerte n'étant branchée sur la table des détections, la chaîne se termine par un écran
et le délai perçu dépend de la présence de l'analyste devant le tableau de bord --- premier des
trois seuils de \cref{sec:socle-service}.

L'évaluation de l'enrichissement répond au verrou V3 du \cref{ch:etat-art} : les modèles de rapprochement
sémantique sont habituellement évalués sur des textes bien formés, non sur des alertes courtes et
bruitées. Une limite la précède : l'enrichissement n'alimente pas le score d'incident
(\cref{subsec:avancement-enrichissement}), et son effet sur la priorisation n'est donc pas
évaluable. Ce qui est mesurable rigoureusement est la qualité du rattachement d'une alerte à une
technique, comparée à deux méthodes de référence plus simples --- correspondance par mots-clés (M0)
et similarité par pondération de termes sans apprentissage de domaine (M1) --- face à la méthode
retenue (M2) ; protocole, jeu d'évaluation et grille de résultats figurent en
\cref{ann:evaluation-semantique}.

L'enrichissement reste utile sans cette jointure, puisqu'il fournit à l'analyste, depuis chaque
alerte, une technique candidate et la procédure d'incident associée --- seule forme sous laquelle
il l'atteint, et que la \cref{fig:capture-K33} montre à l'écran. En l'absence de mesure
comparative, le rapport se garde en revanche d'affirmer ce rattachement assez précis pour être
automatisé dans le calcul de priorité --- l'objectif O5 avait par avance admis qu'un résultat
négatif serait rapporté tel quel. Une validation distincte, déjà exécutée, mérite d'être signalée
sans s'y substituer : un export quantifié du modèle a été rejeté après un contrôle d'accord de
classement mené le 1\textsuperscript{er} août 2026 --- 0~correspondance de premier rang sur
5~requêtes face au modèle de référence --- et le format non quantifié a été conservé en recette, au
prix d'une empreinte mémoire plus élevée. Ce petit échantillon invalide cet export précis ; il ne
démontre ni la qualité générale du modèle retenu ni sa supériorité sur M0, M1 ou M2.

\begin{figure}[!htbp]
\centering
\IfFileExists{figures/K33.png}{%
  \setlength{\fboxrule}{0.4pt}\setlength{\fboxsep}{0pt}%
  \fcolorbox{black!30}{white}{\includegraphics[width=0.72\textwidth]{figures/K33.png}}%
}{%
  \begin{minipage}{0.72\textwidth}
  \begin{extraitconf}[breakable=false, fontupper=\footnotesize,
      title={Cadre de capture K33 : spécification de prise}]
  \textbf{Vue :} fiche d'incident du tableau de bord, une entité agrégée sélectionnée depuis
  la page \emph{Incidents}, la tâche d'enrichissement ayant tourné sur la fenêtre courante.
  \textbf{Date de prise :} relevée au rejeu, et portée dans la légende.\par\smallskip
  \textbf{Ce que le cadre établit :} ce que l'enrichissement rend à l'analyste --- la technique
  candidate et son score de similarité affichés à côté de l'alerte, la jauge de score d'incident
  et la chronologie de l'entité. Il établit du même coup la borne rapportée au même endroit : le
  score de la jauge ne dépend pas du candidat sémantique affiché.\par\smallskip
  \textbf{Masquage :} identifiant de projet, adresses réelles.
  \end{extraitconf}
  \end{minipage}%
}
\caption[Fiche d'incident : le candidat sémantique rendu à l'analyste]{Fiche d'incident du
tableau de bord. À y lire : la technique candidate et sa similarité affichées à côté de
l'alerte --- c'est la seule forme sous laquelle l'enrichissement atteint l'analyste --- et
la jauge de score, calculée sans ce candidat, qui rend visible l'écart d'architecture
énoncé au même endroit.}
\label{fig:capture-K33}
\end{figure}

La boucle reliant les vulnérabilités livrées aux techniques observées se mesure par une question
unique : l'ordre de priorité qu'elle produit diffère-t-il de l'ordre par gravité théorique seule\,?
Y répondre suppose un rattachement des vulnérabilités aux techniques, dont le champ est laissé vide
par le script de chargement. F6 est donc un prototype dont la chaîne fonctionne de bout en bout ---
les constats remontent réellement dans l'entrepôt et s'affichent --- mais dont la priorisation n'est
pas validée : \textbf{la boucle transporte la donnée, elle n'en tire pas encore d'ordre}.

\section{Performance, résilience et coût}
\label{sec:performance}

\subsection{Objectifs de service}

\textbf{Quatre objectifs de service} sont configurés sur trente jours glissants : disponibilité et
latence, pour la plateforme et pour l'application hébergée \cite{gcp_architecture_framework}. La
disponibilité est fixée à 99\,\% --- budget d'erreur de 7~h~12 sur 30~jours --- et la latence à
95\,\% des requêtes sous 1~s. Deux décisions de conception les accompagnent, justifiées dans le
code : les erreurs clientes sont exclues du numérateur de disponibilité, car compter un refus
d'authentification comme une indisponibilité ferait chuter l'objectif précisément quand la
plateforme fait son travail ; et le seuil de latence a été relâché de 800~ms à 1~s en conséquence du
choix de réduire à zéro le nombre d'instances actives hors trafic --- \textbf{un seuil que l'on sait
déjà intenable ne mesure rien, il apprend seulement à ignorer l'alerte}. Ce sont des objectifs
configurés, non atteints, faute de trafic réel (\cref{sec:cloture}) ; la seule observation
disponible porte sur 48~heures relevées le 07/08/2026, où la latence au
95\textsuperscript{e}~centile de l'API est restée entre 9,5~ms et 1004~ms, un seul point au-dessus
de 800~ms --- fenêtre trop courte pour valider l'objectif, ordre de grandeur tenu en recette.

\textbf{Observabilité de sécurité et observabilité d'exploitation.} Les deux ne sont pas au même
stade et les confondre serait surestimer le résultat. L'\emph{observabilité de sécurité} est
réalisée et démontrée : collecte filtrée F4, sept règles actives, table d'indicateurs agrégés
(volume, erreurs, latence, blocages), et un scénario complet vérifié le 19/08/2026.
L'\emph{observabilité d'exploitation} est configurée sans être éprouvée : \textbf{dix-sept
politiques d'alerte réparties en onze familles} --- des taux d'erreur et de latence à l'arrêt de
l'ingestion et à la dérive de l'isolation entre locataires --- déclarées en code et acheminées vers
un canal de notification, complétées par \textbf{deux sondes de disponibilité} interrogeant les
points de santé de l'API et de l'application hébergée, seuls signaux de vivacité indépendants du
trafic.

Trois de ces politiques ont depuis été provoquées et portent leur mesure datée au
\cref{tab:synthese-campagne} : absence de trafic (T12), retard d'enrichissement ouvert sans avoir
été provoqué (T15), budget de facturation (T19). Deux bornes subsistent et portent le même défaut :
les quatorze autres politiques n'ont jamais été déclenchées, et les quatre objectifs de service
restent mesurés sans être actionnés, faute de politique d'alerte sur la consommation de leur budget
d'erreur.

Le démarrage à froid du composant d'encodage sémantique porte une mesure réelle de début
août~2026 (\cref{graph:coldstart}) : moyenne $\approx$~27~s, pic $\approx$~94~s, très
au-dessus de l'estimation de conception, de l'ordre de deux à quatre secondes --- conservée plutôt
que corrigée en silence, pour que l'écart au mesuré reste lisible (incident et traitement en
\cref{subsec:incidents}). Le tableau de bord présente un profil plus favorable : environ 10,5~s le
07/08/2026, dominées par l'initialisation de la plateforme d'exécution et non par le code.

\begin{figure}[!htbp]
\centering
\begin{tikzpicture}[y=0.0185cm, x=0.82cm, font=\footnotesize]
% Aire de trace bornee juste apres la seconde barre ; graduation reguliere de 25 en 25.
\draw[thick] (-0.3,0) -- (4.9,0);
\draw[thick] (-0.3,0) -- (-0.3,160);
\node[rotate=90, anchor=south, font=\scriptsize] at (-1.25,80) {Secondes};
\foreach \y in {0,25,50,75,100,125,150}{
  \draw (-0.45,\y) -- (-0.3,\y);
  \node[anchor=east, font=\scriptsize] at (-0.5,\y) {\y};
  \draw[black!20] (-0.3,\y) -- (4.9,\y);
}
\fill[black!12] (0.4,0) rectangle (1.8,27); \draw[thick] (0.4,0) rectangle (1.8,27);
\node[anchor=south, font=\scriptsize\bfseries] at (1.1,27) {27~s};
\node[anchor=north, font=\scriptsize, align=center] at (1.1,-4) {Moyenne mesurée};
\fill[black!45] (2.8,0) rectangle (4.2,94); \draw[thick] (2.8,0) rectangle (4.2,94);
\node[anchor=south, font=\scriptsize\bfseries] at (3.5,94) {94~s};
\node[anchor=north, font=\scriptsize, align=center] at (3.5,-4) {Pic observé (p99)};
\draw[dashed, very thick] (-0.3,150) -- (4.9,150);
\node[anchor=south west, font=\scriptsize] at (-0.25,152) {budget de la sonde porté à 150~s (07/08/2026)};
\end{tikzpicture}
\caption[Démarrage à froid du service d'encodage]{Démarrage à froid du service d'encodage : moyenne et pic mesurés début août 2026,
comparés au budget de la sonde de disponibilité.}
\label{graph:coldstart}
\end{figure}

\subsection{Reproductibilité et résilience}
\label{sec:resilience}

La reconstruction complète d'un environnement depuis le seul dépôt répond au constat
d'infrastructure non reproductible du \cref{ch:cadre-existant}, et deux questions s'y distinguent. La
\textbf{capacité} est acquise : la recette a été construite, puis reconstruite, à partir du seul
code d'infrastructure --- c'est ce qui fonde la réponse à O2. Sa \textbf{mesure}, durée de bout en
bout en une exécution continue, relève du plan de validation continue à cadence semestrielle et
sera exécutée en direct en soutenance. \emph{Réserve} : le fichier de variables de l'environnement
n'est pas versionné, de sorte qu'une reconstruction depuis un poste neuf ne reproduit pas encore
l'environnement à l'identique --- le versionner, ce fichier ne contenant aucun secret, est le
correctif de reproductibilité le plus rentable identifié et le prérequis de la mesure. La
restauration de la base, elle, a été exécutée et chronométrée (\cref{tab:restore-db}).

\begin{table}[!htbp]
\caption[Restauration de la base, test réel du 03/08/2026]{Restauration de la base de données --- test réel exécuté le 03/08/2026}
\label{tab:restore-db}
\centering
\footnotesize
\setlength{\tabcolsep}{3pt}
\begin{tabular}{|P{4.6cm}|P{10.9cm}|}
\hline
\textbf{Grandeur} & \textbf{Valeur mesurée} \\
\hline
RTO (durée de l'opération de restauration) & 32~min~45~s (opération de clonage, de
\hms{11:16:40} à \hms{11:49:26} UTC, instance d'entrée de gamme, 10~Go) \\ \hline
RPO (perte de données au point demandé) & 0 \\ \hline
Résultat au point demandé & 314~lignes attendues, 314~restaurées, aucune manquante ni en trop ;
les trois écritures postérieures au point de restauration correctement absentes \\ \hline
Dispositif de sauvegarde & 7~jours de journaux de transactions conservés en continu ;
7~sauvegardes complètes, quotidiennes à \hms{02:00} UTC \\ \hline
Portée de la mesure & Configuration zonale, antérieure au passage en haute disponibilité
régionale du 08/08/2026 ; mesure à revalider sur la configuration actuelle \\
\hline
\end{tabular}
\end{table}

Le RTO déclaré ne couvre que le clonage technique : détection, décision et bascule applicative, non
chronométrées séparément, portent l'estimation de bout en bout entre 45~min et 1~h. Une précision
d'exploitation a des conséquences contractuelles : la même instance héberge les deux bases, si bien
qu'un clone de restauration les restaure ensemble. Le retour arrière applicatif par révision est
mesuré le 07/08/2026 sur l'API en recette : \textbf{11,6~s} vers la révision précédente,
\textbf{16,5~s} au retour, l'écart tenant à un démarrage à froid --- mesure non répétée, ni le
nombre de séries ni le délai de validation humaine n'étant documentés, donc ordre de grandeur et
non objectif de rétablissement.

Le mécanisme de détection de dérive \cite{hashicorp_drift} (\cref{subsec:derive}) a été sollicité
en conditions réelles début août~2026, lors de l'incident de chiffrement géré par le client. Cet
épisode ne vaut pas T16 corrigé, une ressource créée hors état échappant à la seule planification.
Le seul contrôle de dérive \emph{planifié} est d'une autre nature --- la tâche quotidienne qui
vérifie la persistance du durcissement de l'isolation entre bases (\cref{sec:locataire}) --- et
c'est un aveu d'architecture autant qu'un contrôle : une propriété qu'il faut surveiller
quotidiennement n'est pas déclarative.

\subsection{Coût réel}
\label{sec:cout}

Aucune ventilation mensuelle vérifiable par poste n'existe dans les sources, et aucun relevé de
facturation ventilé n'est donc publié. Un fait reste sourcé : le document d'architecture du projet
situe le coût au niveau des paliers gratuits, à quelques dizaines d'euros par mois, désigne
l'ingestion de journaux comme poste dominant et confirme l'absence de processeur graphique et de
base vectorielle facturée séparément --- constat interne daté, non relevé de facturation. Trois
constats datés le complètent (\cref{tab:cout-anomalies}).

\begin{table}[!htbp]
\caption[Constats de coût datés]{Constats de coût datés, statut à la date de \dateref}
\label{tab:cout-anomalies}
\centering
\footnotesize
\setlength{\tabcolsep}{3pt}
\begin{tabular}{|P{3.0cm}|P{5.2cm}|P{7.1cm}|}
\hline
\textbf{Poste et statut} & \textbf{Constat daté} & \textbf{Direction de remédiation} \\
\hline
Services d'exécution --- \textbf{corrigé} & Instance maintenue active en permanence, donc coût
continu sans contrepartie en recette ; réglage hérité jamais déclaré & Nombre minimal d'instances
ramené à zéro hors trafic, arbitrage coût contre latence déclaré \\ \hline
Requêtes planifiées de l'entrepôt --- \textbf{ouvert} & $\approx$~78~Go/j facturés pour 24~Mo de
données : minimum forfaitaire par table balayée $\times$ exécutions quotidiennes & Réduire le
nombre d'exécutions facturées plutôt que le volume : fusionner les sept requêtes en une seule
passe, restreindre le balayage aux partitions de la fenêtre utile, et n'espacer la cadence qu'en
dernier recours, une cadence plus lente dégradant le délai de détection \\ \hline
Haute disponibilité base de données --- \textbf{appliqué} & Bascule régionale appliquée et
confirmée en direct le 08/08/2026 ; basculement de zone jamais provoqué & Provoquer un
basculement de zone en fenêtre maîtrisée et le chronométrer, seul moyen d'établir le gain payé \\
\hline
\end{tabular}
\end{table}

Deux conséquences que le tableau ne porte pas. Le volume facturé venant du minimum forfaitaire et
non des données, \emph{une réduction de rétention n'aurait ici aucun effet sur la facture}. Et le
dimensionnement est incohérent, ce qui se dit franchement : l'instance retenue est de gamme
partagée, déconseillée en production, tout en étant facturée au tarif de la haute disponibilité
régionale --- appliquée en 11~min~2~s le 08/08/2026, dont la preuve \emph{dynamique} suppose un
basculement de zone provoqué, porté par le plan de validation continue. Aucune des trois directions
n'est appliquée à la date de \dateref{} : elles sont énoncées comme suite à donner, non comme
correction acquise.

Faute de volume représentatif et de grille tarifaire datée, l'alternative commerciale étudiée au
\cref{ch:etat-art} n'est pas chiffrée. La comparaison \emph{qualitative} reste instructive : les deux
options partagent la même variable dominante, l'ingestion, que la solution commerciale facture avec
un abonnement plancher et dont le socle se passe, aux volumes de la recette, dans les paliers
gratuits du fournisseur. \textbf{L'écart de structure compte plus que l'écart de montant} : on
achète dans un cas un service exploité par un tiers, dans l'autre du temps d'ingénieur et une dette
de maintenance --- dont le plafond de règles traduisibles par une personne est la contrepartie.

\section{Démonstration de bout en bout et conformité aux référentiels}
\label{sec:demonstration}

Le 19/08/2026, un scénario d'attaque complet a été exécuté en recette : il combine les preuves de
cinq domaines --- identité, périmètre, cloisonnement des données, supervision et chaîne de
livraison --- en une démonstration continue que la \cref{fig:scenario-19-08} déroule de l'attaque à
l'alerte qualifiée ; les treize requêtes mêlaient injections SQL, script inter-sites, inclusion de
fichier local et accès à des fichiers sensibles. Cette observation ponctuelle ne mesure pas le taux
de faux positifs, qui n'a pas été établi.

\begin{figure}[!htbp]
\centering
\begin{tikzpicture}[
  jalon/.style={draw, thick, rounded corners=2pt, fill=black!4, text width=3.0cm,
    minimum height=1.85cm, align=center, font=\footnotesize, inner sep=3pt},
  pastille/.style={circle, draw, thick, fill=black, inner sep=1.5pt},
  h/.style={font=\scriptsize\bfseries, align=center}
]
\draw[-{Latex[length=2.2mm]}, thick] (-1.5,0) -- (13.4,0);
\node[font=\scriptsize, anchor=south west] at (-1.5,0.10) {heure UTC};   % place au-dessus de l'axe : sous l'axe il chevauchait le premier horodatage
\foreach \x in {0,4,8,12}{\node[pastille] at (\x,0) {}; \draw[thick] (\x,0) -- (\x,0.7);}
\node[jalon, anchor=south] at (0,0.7)
  {13~requêtes SQLi, XSS et LFI envoyées vers l'application hébergée};
\node[jalon, anchor=south] at (4,0.7)
  {13 réponses~403 au périmètre ; 1~requête légitime en~200};
\node[jalon, anchor=south] at (8,0.7)
  {13~lignes de journaux portant le service du locataire};
\node[jalon, anchor=south] at (12,0.7)
  {Alerte R2 : « 13~requêtes bloquées en 15~min », TA0040 / T1498};
\node[h, anchor=north] at (0,-0.2)    {\hms{16:41:xx}};
\node[h, anchor=north] at (4,-0.2)    {\hms{16:41:5x}};
\node[h, anchor=north] at (8,-0.2)    {quelques secondes\\ (non horodaté)};
\node[h, anchor=north] at (12,-0.2)   {\hms{16:56:08}\\ ($\approx$~15~min)};
\node[anchor=north west, text width=15.5cm, align=left, font=\scriptsize, inner sep=0pt]
  at (-1.5,-1.5)
  {\emph{Échelle non linéaire.} L'étiquetage TA0040 / T1498 est celui produit par la règle ;
   une qualification par T1190 serait plus plausible pour ces charges et doit être revue par
   l'analyste. L'enrichissement disponible datait de \hms{12:16:15} : il n'enrichit pas cette
   alerte.};
\end{tikzpicture}
\caption[Scénario de bout en bout du 19/08/2026]{Scénario de bout en bout du 19/08/2026 : de l'attaque à l'alerte qualifiée}
\label{fig:scenario-19-08}
\end{figure}

L'alerte R2 du 19/08/2026 à \hms{16:56:08} UTC est le dernier maillon de la chaîne : elle compte
treize requêtes bloquées sur sa fenêtre de 15~min et les attribue au locataire --- attribution qui
est l'exception plutôt que la règle (\cref{sec:socle-service}). Deux contrôles complémentaires ont
été relancés en direct le même jour : le contrôle d'isolation entre bases a produit \textbf{six
vérifications} toutes en échec comme attendu --- une connexion croisée depuis l'utilisateur d'un
locataire vers la base de l'autre, réellement tentée et rejetée --- et le compte de service du
second locataire dispose exactement de deux rôles, sans droit sur l'entrepôt. Le dernier
rattachement sémantique disponible, daté de \hms{12:16:15} UTC, présente une similarité de 0,698
avec T1556.003 : précédant l'attaque de plus de quatre heures, il atteste le fonctionnement
antérieur du composant, pas l'enrichissement de l'alerte R2.

Le premier maillon de cette chaîne --- les treize refus eux-mêmes --- est le seul dont la preuve
tienne dans une seule vue : la \cref{fig:capture-K14} le restitue.

\begin{figure}[!htbp]
\centering
\IfFileExists{figures/K14.png}{%
  \setlength{\fboxrule}{0.4pt}\setlength{\fboxsep}{0pt}%
  \fcolorbox{black!30}{white}{\includegraphics[width=0.86\textwidth]{figures/K14.png}}%
}{%
  \begin{minipage}{0.86\textwidth}
  \begin{extraitconf}[breakable=false, fontupper=\footnotesize,
      title={Cadre de capture K14 : spécification de prise}]
  \textbf{Vue :} collecteur de journaux, requête filtrée sur les journaux du répartiteur pour la fenêtre du 19/08/2026 autour de \hms{16:41} UTC, colonnes horodatage, code de réponse et nom de la règle de filtrage appliquée.\par\smallskip
  \textbf{Date de prise :} 19/08/2026 --- la fenêtre est datée par le scénario lui-même ; la capture se reprend depuis le collecteur tant que la rétention le permet, sans rejeu.\par\smallskip
  \textbf{Contenu probant :} les \textbf{treize} requêtes de la \cref{fig:scenario-19-08} refusées en \textbf{403} au périmètre, horodatées dans la fenêtre annoncée, avec le nom de la règle qui a prononcé chaque refus. C'est ce décompte, et lui seul, que l'alerte R2 de \hms{16:56:08} agrège quinze minutes plus tard. \textbf{Le cadre ne vaut pas preuve de la couverture du filtrage} : le motif brut de traversée de chemin, qui a reçu une redirection 302 le 23/08/2026, n'appartient pas à ce scénario (\cref{sec:non-conformite-t1}).\par\smallskip
  \textbf{Masquage} (rectangle opaque, jamais de flou) \textbf{:} numéro de projet cloud, adresses source réelles des requêtes.
  \end{extraitconf}
  \end{minipage}%
}
\caption[Les treize refus au périmètre du 19/08/2026 dans le collecteur de journaux]{Les treize requêtes du scénario du 19/08/2026 refusées en 403 au périmètre, telles qu'elles figurent dans le collecteur de journaux. Restitution du premier maillon de la \cref{fig:scenario-19-08} ; elle illustre une preuve nommée et ne lui tient pas lieu.}
\label{fig:capture-K14}
\end{figure}

\subsection{Correspondance avec les preuves de soutenance, et bornes du scénario}

Le scénario ne vaut comme démonstration que s'il se rattache aux résultats du rapport, et les cinq
preuves attendues en soutenance s'y retrouvent chacune : le \textbf{périmètre} par les treize
requêtes bloquées en~403 (\cref{sec:bout-en-bout}, test T1) ; l'\textbf{identité} par les deux rôles exacts et
l'usurpation refusée (\cref{sec:plan-identite}, test T5) ; le \textbf{cloisonnement des données} par les six
vérifications d'isolation sur six (\cref{sec:locataire}, test T11) ; la \textbf{supervision} par l'alerte R2 de
\hms{16:56:08} UTC, dont la qualification doit être revue par l'analyste
(\cref{sec:chaine-detection}) ; et la
\textbf{chaîne de livraison} par l'exécution réelle de l'analyse statique, en échec sur 74
constats (\cref{sec:chaine-livraison}, incident de la porte inerte ; test T8). Trois points restent hors du périmètre du
scénario, tous rattachés à des protocoles de la campagne : la ségrégation réseau entre locataires,
écart documenté en \cref{sec:socle-service} ; le motif brut de traversée de chemin
(\cref{sec:non-conformite-t1}) ; et les constats d'analyse statique --- 74 constats et une tâche
en échec le 19/08, dont la remédiation complète n'est pas démontrée.

\subsection{Conformité aux référentiels et maturité Zero Trust}
\label{sec:maturite}

Le positionnement par rapport aux principes Zero Trust \cite{cisa_zt_maturity_model} et au
référentiel de configuration est celui du \cref{ch:conception} : ni certification, ni audit exhaustif, et
aucun taux global sans grille contrôle par contrôle. Trois des onze écarts du
\cref{ch:conception} touchent cette campagne, et les trois sont clos : journalisation des
refus réseau (É1) et identité partagée du composant d'encodage (É3), traitées en
\cref{sec:preuve-t11} ; deux sous-réseaux inutilisés (É2), retirés le 24/08/2026
(\cref{subsec:migration-reseau}). Le mémoire prononce en revanche lui-même son niveau de maturité
pilier par pilier, avec une justification chiffrée et un résultat \emph{asymétrique}
(\cref{tab:ztmm}) : \textbf{un profil uniformément élevé serait un profil complaisant}, les piliers
ne progressant jamais au même rythme parce qu'ils ne coûtent pas le même prix.

{\footnotesize
\setlength{\tabcolsep}{3pt}
\begin{longtable}{|P{2.5cm}|P{1.7cm}|P{11.1cm}|}
\caption[Auto-évaluation de maturité Zero Trust]{Auto-évaluation sur le modèle de maturité Zero Trust, à la date de \dateref}
\label{tab:ztmm} \\
\hline
\textbf{Pilier} & \textbf{Niveau} & \textbf{Justification chiffrée et borne du niveau} \\
\hline
\endfirsthead
\hline
\textbf{Pilier} & \textbf{Niveau} & \textbf{Justification chiffrée et borne du niveau} \\
\hline
\endhead
Identités & \textbf{Avancé} & Sept identités de service distinctes, aucun rôle primitif sur 53
fichiers de description d'infrastructure, aucune clé de compte de service, fédération contrainte
sur le dépôt \emph{et} la branche, jetons courts, second facteur temporel sur le tableau de bord.
\emph{Borne} : second facteur résistant à l'hameçonnage, codes de secours, garantie préventive par
politique d'organisation \\ \hline
Terminaux & \textbf{Traditionnel} & Aucune notion d'appareil dans toute la chaîne : ni inventaire,
ni certificat client, ni signal de posture ; le seul attribut de contexte est le pays d'origine.
\textbf{C'est le pilier vide du modèle}, par absence de périmètre administratif et non par
oubli \\ \hline
Réseaux & \textbf{Initial à avancé} & Nord-sud solide : filtrage applicatif appliqué, transport
chiffré imposé, entrée des services restreinte au répartiteur, base sans adresse publique.
Est-ouest refondu du 11 au 25/08/2026 (décision D15, \cref{subsec:reseau-cible}) : un sous-réseau
par locataire, refus par défaut en sortie posé le 23/08/2026, cinq autorisations seulement, chacune
ciblée par une identité de service et par un port, aucune plage source ; la règle interne
permissive a été supprimée le 24/08/2026 et le connecteur partagé le 15/08/2026. Un refus croisé
entre locataires a été réellement opposé et journalisé le 24/08/2026 (T20, conforme).
\emph{Borne} : deux charges de travail détachées du sous-réseau conservent une sortie managée hors
de portée du pare-feu (T14, non conforme, 25/08/2026) ; aucune règle de détection ne lit les refus
journalisés ; l'instance de base de données reste jointe par les deux locataires \\ \hline
Applications et charges de travail & \textbf{Avancé} & Entrée par le répartiteur seule sur tous
les services, composant d'encodage en accès interne strict, contrôle des rôles vérifié en direct
sur les sept routes serveur, en-têtes de sécurité durcis, conteneurs sans privilège, porte
d'analyse d'image bloquante et \emph{éprouvée} (T8). \emph{Borne} : documentation d'interface
publique, droit d'invocation anonyme adossé au seul contrôle de périmètre \\ \hline
Données & \textbf{Initial à avancé} & Chiffrement par clé gérée sur l'entrepôt, les secrets et le
stockage d'objets ; séparation d'écriture par table \emph{vérifiée par test automatisé} (T11) ;
rétention de 90~jours sur les journaux bruts ; restauration mesurée à 32~min~45~s avec perte
nulle. \emph{Borne} : chiffrement non rétroactif sur les tables de preuves antérieures, aucune
protection contre la suppression, et trois seulement des dix tables de l'entrepôt portent une clé
d'application, une seule la portant renseignée sur la totalité de ses lignes
(\cref{tab:modele-donnees}) \\ \hline
\textbf{Visibilité et analyse} & \textbf{Avancé} & Quatre acheminements de journaux, trafic du
répartiteur collecté sans filtre, sept règles de détection avec étiquetage de technique,
enrichissement sémantique, et un tableau de bord qui affiche aussi les \emph{non}-détections \\ \hline
\textbf{Automatisation et orchestration} & \textbf{Initial à avancé} & Infrastructure décrite en
code à l'exception du nom de domaine ; chaîne de livraison fédérée sans clé ; requêtes planifiées
et contrôle de dérive quotidien de l'isolation. \emph{Borne} : infrastructure appliquée
manuellement, sans planification en chaîne ni détection de dérive programmée --- le maillon le
moins mature (T16, T18) \\ \hline
\textbf{Gouvernance} & \textbf{Initial} & Dix-sept décisions d'architecture datées et un registre
d'écarts tenu sans complaisance. \emph{Borne} : aucune politique d'organisation, aucun processus
de révision périodique des règles de filtrage par pays \\
\hline
\end{longtable}
}

La formulation à défendre tient en une phrase : le socle met en \oe uvre un Zero Trust
\emph{nord-sud} mature --- identité, périmètre, charges de travail --- et un est-ouest gouverné par
l'identité \emph{au plan réseau} depuis le 25/08/2026, mais non au plan du stockage, où l'instance
de base de données et l'entrepôt restent communs aux deux locataires. Ce qui reste périmétrique
l'est par contrainte de périmètre, tracée en décision : sans organisation au sens du fournisseur,
ni politique d'organisation, ni périmètre de service, ni accès conditionnel géré ne sont
disponibles. Le seuil de rupture de ce compromis est l'accueil d'un troisième locataire.

\section{Ce que le filtrage applicatif ne protège pas}
\label{sec:limites-filtrage}

Un chapitre de validation qui n'énoncerait que ce qui fonctionne validerait mal. Deux limites du
socle sont structurelles : elles ne se corrigent pas par un réglage, et elles délimitent exactement
ce qu'un hébergeur peut promettre.

\subsection{Invisibilité d'une faille de logique métier}

Le risque le plus concret porté par l'application hébergée est le suivant : le signalement d'une
phrase du corpus la désactive aussitôt, sans quorum. Un contributeur authentifié, agissant dans
son quota, peut donc désactiver le corpus phrase par phrase --- l'actif économique de
l'application --- en émettant des requêtes parfaitement bien formées.

Ce que voit le socle se résume en un mot : rien. Le filtrage applicatif observe du trafic
légitime --- méthode autorisée, chemin fonctionnel, charge utile sans motif d'attaque --- et aucune
des sept règles de détection ne se déclenche : R1 compte les refus d'authentification, or l'utilisateur est
authentifié ; R2 compte les blocages du filtrage, or il n'y en a aucun ; R3 cherche des séquences
de traversée de chemin, absentes ; R4 cherche des agents utilisateurs scriptés, et le service de
l'application hébergée en est de toute façon explicitement exclu ; R5 surveille la latence,
inchangée ; R6 cherche des motifs d'injection, absents ; R7 surveille l'accès à des fichiers
sensibles, non sollicités.

La conclusion s'écrit telle quelle : \textbf{un filtrage applicatif ne corrige pas une faille de
logique métier, et une chaîne de détection fondée sur des signaux de transport ne la voit pas}. Le
socle n'observe ni les actions fonctionnelles de l'application, ni son modèle d'autorisation, ni
ses invariants métier. Le correctif --- un quorum avant désactivation --- est une migration
applicative et appartient à l'éditeur : c'est la justification concrète de la frontière de
responsabilité de \cref{tab:frontiere-resp}, conséquence de ce qu'un dispositif de transport peut
observer.

\subsection{Portée de la limitation de débit au bord}

La seconde limite est du même ordre. La limitation de débit posée au périmètre est indexée sur
l'adresse source : efficace contre une salve provenant d'une adresse unique, ce que T2 a vérifié,
elle est sans effet contre un adversaire distribué qui répartit le même volume sur assez d'adresses
pour que chacune reste sous le seuil.

Le symétrique est plus gênant encore, et il a été rencontré. Là où de nombreux abonnés partagent un
petit nombre d'adresses publiques de sortie --- situation courante sur les réseaux mobiles du pays
visé par l'application hébergée, et anticipée par son éditeur --- un seuil par adresse frappe des
utilisateurs légitimes : le compteur agrège des personnes différentes. C'est la généralisation de
l'incident rapporté en \cref{subsec:authent-dashboard}, dont la formulation est que \emph{le
contrôle mesure la mauvaise entité}. Tant que la limitation reste au bord, la resserrer bloque
davantage d'utilisateurs légitimes sans gêner l'adversaire distribué, et la relâcher fait
l'inverse : la sortie du compromis n'est pas un meilleur seuil, c'est une limitation portée par
l'identité, donc côté application, ou indexée au bord sur un jeton de session --- ce qui suppose
que le bord sache lire cette session.

\section{Du socle au service}
\label{sec:socle-service}

Le \cref{ch:cadre-existant} posait le socle comme un produit potentiel, non comme un exercice. Cette section en
tire le bilan avec les mêmes exigences de preuve que le reste du chapitre : ce que le socle apporte
sans modification de l'application hébergée, ce qu'il n'apporte pas, ce que coûte l'accueil d'une
application supplémentaire, et à quelles conditions l'ensemble deviendrait contractualisable.

\subsection{Apport du socle et charge de l'éditeur}

L'accueil de l'application pilote a servi de test grandeur nature de la proposition de valeur.
Sont apportés automatiquement, sans une ligne de code applicatif : le certificat de transport et sa
redirection, une adresse d'entrée unique, le filtrage applicatif et sa limitation de débit
anti-force-brute avec bannissement, la journalisation centralisée et sa normalisation, les sept
règles de détection avec leur étiquetage de technique, l'enrichissement sémantique, la restitution
dans le tableau de bord, les alertes d'exploitation, une identité de service dédiée à moindre
privilège, un secret par usage chiffré par clé gérée, une base et un espace de stockage dédiés, et
une chaîne de livraison avec détection de secrets, analyse statique et double analyse d'image
bloquante.

Une preuve précise vaut mieux qu'un argument général. L'application hébergée protégeait ses accès
d'administration par une liste d'adresses autorisées, lue depuis un en-tête de proxy propre à son
hébergeur précédent ; derrière le répartiteur du socle, cet en-tête devient forgeable par le client,
et le contrôle d'administration se serait effondré silencieusement au moment de la migration. Il a
été corrigé pendant l'accueil, par un arbitrage explicite de l'en-tête de confiance selon
l'hébergeur, et le correctif est protégé par un test unitaire érigé en porte bloquante.
\textbf{Changer d'hôte peut casser un contrôle de sécurité, et seul un socle qui l'a anticipé le
rattrape.}

\textbf{Ce qui incombe à l'éditeur.} La frontière est établie au \cref{ch:realisation}
(\cref{tab:frontiere-resp}) ; ce qu'elle exclut du périmètre de l'hébergeur tient en six
domaines : authentification et autorisation des utilisateurs finaux --- le second facteur du socle
protège les comptes du tableau de bord, pas les utilisateurs de l'application hébergée ---
validation des entrées, le filtrage ne traitant que les charges connues, sécurité de la logique
métier, structurellement invisible (\cref{sec:limites-filtrage}), chiffrement applicatif des
données sensibles, sauvegardes applicatives, et conformité réglementaire, pour laquelle le socle
fournit des moyens techniques et non une gouvernance. La ségrégation réseau entre applications
hébergées appartient en revanche à l'hébergeur et est couverte au plan réseau depuis la migration
du 11 au 25/08/2026 ; ce qui lui incombe encore est le cloisonnement de la donnée, troisième seuil
énoncé plus bas.

\subsection{Coût d'accueil d'une application supplémentaire}

Le coût d'accueil d'une troisième application est mesuré (\cref{tab:cout-onboarding}) ; la colonne
d'automatisation est la seule qui compte pour un raisonnement de produit, car elle sépare ce qui
passe à l'échelle de ce qui exige un geste humain non versionné, donc non auditable.

\begin{table}[!htbp]
\caption[Coût d'accueil d'une application supplémentaire]{Coût d'accueil d'une application hébergée supplémentaire, à la date de \dateref}
\label{tab:cout-onboarding}
\centering
\footnotesize
\setlength{\tabcolsep}{3pt}
\begin{tabular}{|P{3.4cm}|P{7.0cm}|P{4.9cm}|}
\hline
\textbf{Nature} & \textbf{Détail} & \textbf{Automatisé\,?} \\
\hline
Un fichier de description ($\approx$~120~lignes) & Sur le modèle de celui du locataire existant &
Recopie, module réutilisable \\ \hline
Cinq variables d'environnement & À mettre à jour ensemble ; un oubli rend l'application absente
de la supervision \emph{sans aucune erreur visible} & Manuel, et le fichier n'est pas versionné \\ \hline
Un flux de livraison ($\approx$~296~lignes) & Duplication intégrale, faute de factorisation &
Manuel \\ \hline
Sept ordres de base de données \emph{pour un locataire supplémentaire} (dix pour l'accueil du
second, qui incluait la fermeture symétrique) & Durcissement de l'isolation entre bases & Hors
description en code, geste appliqué à l'exécution \\ \hline
Une tâche, un ordonnanceur, une alerte & Contrôle quotidien de dérive de l'isolation & À recopier \\ \hline
Deux modifications de code & Table des services du locataire et filtre d'affichage & Manuel,
sans test \\ \hline
Un enregistrement de nom de domaine & À poser \emph{avant} l'application de l'infrastructure &
Manuel, hors périmètre du code \\
\hline
\end{tabular}
\end{table}

C'est la mesure honnête de la distance entre un socle qui \emph{fonctionne} pour deux applications
et un socle qui \emph{s'exploite} pour N.

\subsection{Maturité multi-locataire, axe par axe}

L'évaluation axe par axe (\cref{tab:maturite-multitenant}) est franchement asymétrique, et
c'est ce qui la rend utilisable : elle indique où investir. Le socle a résolu le plan de l'identité
et du déploiement, et, du 11 au 25/08/2026, celui du réseau (D15) ; celui de la donnée reste ouvert
\emph{en exploitation} sans l'être \emph{en conception}, D16 en fixant la cible, le coût et l'ordre
(\cref{subsec:reseau-cible,subsec:isolation-donnees}). Les notes mesurent le système déployé, non
la conception qui le corrige, et se lisent sous les conditions de validité du \cref{tab:limites} :
dispositif mono-opérateur, recette sans trafic réel.

{\footnotesize
\setlength{\tabcolsep}{3pt}
\begin{longtable}{|P{3.2cm}|P{1.0cm}|P{11.1cm}|}
\caption[Maturité multi-locataire, axe par axe]{Maturité multi-locataire, axe par axe, à la date de \dateref}
\label{tab:maturite-multitenant} \\
\hline
\textbf{Axe} & \textbf{Note} & \textbf{Justification chiffrée} \\
\hline
\endfirsthead
\hline
\textbf{Axe} & \textbf{Note} & \textbf{Justification chiffrée} \\
\hline
\endhead
Isolation des identités & \textbf{9/10} & Une identité dédiée par application, aucun droit sur
l'entrepôt de supervision, délégation de la chaîne de livraison ciblée sur cette seule identité,
accès aux secrets accordé secret par secret, aucun rôle primitif nulle part \\ \hline
Isolation des secrets & \textbf{9/10} & Un secret par usage, jamais de portée projet, chiffrement
par clé gérée sur huit secrets \\ \hline
Isolation des données applicatives & \textbf{4/10} & Base et utilisateur dédiés sur une instance
partagée ; l'isolation repose sur six révocations appliquées à l'exécution, hors description en
code, qu'une réconciliation consécutive à un clonage défait sans bruit --- d'où le contrôle
quotidien, qui la détecte sans la prévenir \\ \hline
Ségrégation réseau & \textbf{6/10} & Cloisonnement réseau par locataire acquis et daté au
\cref{tab:ztmm} ; le refus croisé entre locataires a été réellement opposé et journalisé le
24/08/2026, sa ligne retrouvée dans l'entrepôt (T20, conforme). La note s'arrête là pour trois
raisons démontrées, non par prudence : deux charges de travail détachées du sous-réseau conservent
une sortie managée hors de portée du pare-feu (T14, non conforme, 25/08/2026) ; le sélecteur
disponible étant l'étiquette et non le compte de service, la frontière est opposable à une charge
compromise et non à un opérateur capable de déployer ; et l'instance de base de données reste
jointe par les deux locataires \\ \hline
Cloisonnement de la supervision & \textbf{2/10} & Un seul jeu de données pour tous les locataires,
aucune politique d'accès par ligne : le filtre est appliqué par l'API, jamais par le stockage. Côté
tableau de bord, le locataire est porté par un témoin de navigation non signé, absent des
revendications du jeton --- confort de lecture, jamais frontière de sécurité, et la décision
d'architecture le dit \\ \hline
Attribution d'un incident à un locataire & \textbf{5/10} & Fiable à 100\,\% pour R1, dégradée pour
R2, R3 et R6, dont la source est le répartiteur : sur les dix tables de l'entrepôt, \textbf{trois
portent une clé d'application et une seule la porte renseignée sur toutes ses lignes}
(\cref{tab:modele-donnees}), et le champ de service est vide sur 75 à 90\,\% des détections de ces
trois règles. Contre-exemple positif : la détection du 19/08/2026 porte bien le service de
l'application hébergée \\ \hline
Reproductibilité de l'accueil & \textbf{5/10} & Module d'application réutilisable et documenté,
mais \textbf{quatre gestes manuels non versionnés} subsistent
(\cref{tab:cout-onboarding}) \\
\hline
\end{longtable}
}

\textbf{Note d'ensemble : 5/10}, contre 4/10 avant la refonte réseau du 11 au 25/08/2026 : l'axe
réseau est le seul à avoir bougé et il ne rachète ni le stockage ni la supervision, qui bornent la
note. Le socle est prêt à héberger plusieurs applications d'un même éditeur --- ce qu'il fait,
correctement --- non les données de deux clients distincts sous engagement contractuel, leur
isolation y restant une propriété des droits du moteur de base et du code applicatif, non du
stockage.

\subsection{Modèle cible et seuils d'engagement}

Les quatre niveaux d'isolation évalués et le modèle retenu --- mutualiser ce qui est coûteux et
sans donnée client, cloisonner ce qui porte la donnée --- sont arrêtés avec leur coût en
\cref{subsec:isolation-donnees}. Ce que la mesure ajoute est le motif de proportion : ce modèle
tranche exactement là où passe la frontière de maturité relevée ci-dessus, pour un surcoût de $+40$
à $+60\,\%$ et un effort \emph{incrémental}. Une limite l'accompagne : tant qu'il n'existe qu'un
seul projet au sens du fournisseur, un droit mal attribué traverse d'un coup toutes les frontières
logiques --- la séparation par projet reste la seule frontière garantie nativement, et lever la
contrainte administrative est le prérequis de toute commercialisation.

\medskip
\noindent\textbf{Trois seuils avant engagement contractuel.} Trois seuils --- exactement trois ---
séparent ces deux états, énoncés dans l'ordre où ils doivent être franchis, celui du rapport entre
l'effet et l'effort.

\textbf{Premier seuil : alerter réellement sur les détections de sécurité.} Le socle exploite deux
chaînes d'observabilité disjointes --- une chaîne d'exploitation qui alerte, une chaîne de sécurité
qui n'alerte pas : une détection de sévérité critique n'envoie aucune notification, elle attend
qu'un analyste ouvre le tableau de bord (\cref{fig:capture-K34}). \textbf{La chaîne se termine par un écran, elle doit se
terminer par une notification} ; c'est la lacune la plus lourde du dispositif, et sa remédiation la
plus légère des trois.

\begin{figure}[!htbp]
\centering
\IfFileExists{figures/K34.png}{%
  \setlength{\fboxrule}{0.4pt}\setlength{\fboxsep}{0pt}%
  \fcolorbox{black!30}{white}{\includegraphics[width=0.72\textwidth]{figures/K34.png}}%
}{%
  \begin{minipage}{0.72\textwidth}
  \begin{extraitconf}[breakable=false, fontupper=\footnotesize,
      title={Cadre de capture K34 : spécification de prise}]
  \textbf{Vue :} page \emph{Vue d'ensemble} du tableau de bord, cadrée sur le bandeau supérieur
  et la bande des cinq indicateurs, un locataire sélectionné et au moins une détection de gravité
  élevée dans la répartition. \textbf{Date de prise :} relevée sur la révision de recette en
  service.\par\smallskip
  \textbf{Ce que le cadre établit :} le premier des trois seuils qui séparent un socle d'un
  service --- la chaîne de sécurité se termine par un écran. Aucune file d'incidents à traiter,
  aucun accusé de prise en charge, aucun canal de notification : cette absence, au milieu de
  compteurs vifs, est l'objet même du cadre, et c'est la raison pour laquelle il est ici plutôt
  qu'une phrase.\par\smallskip
  \textbf{Réserve :} les compteurs de cette vue incluent le trafic d'auto-consultation
  (\cref{tab:limites}) ; la preuve porte sur leur environnement d'interface, non sur leur niveau.
  \textbf{Masquage :} identifiant de projet, courriel du compte connecté.
  \end{extraitconf}
  \end{minipage}%
}
\caption[Vue d'ensemble : la chaîne de sécurité se termine par un écran]{Vue d'ensemble du
tableau de bord. À y lire ce qui \emph{n'y est pas} : ni file d'incidents à traiter, ni
accusé de prise en charge, ni départ de notification. La chaîne de détection se termine par
cet écran, et le délai perçu par un analyste dépend donc de sa présence devant lui.}
\label{fig:capture-K34}
\end{figure}

\textbf{Deuxième seuil : verrouiller la rétention des journaux d'audit.} Les journaux d'accès aux
données sont produits mais ne sont acheminés vers aucune destination dédiée ; ils restent dans un
espace à rétention de trente jours, modifiable. C'est la brique de non-répudiation manquante, et
le contexte la rend impérative : la table des détections a réellement été vidée par une
suppression administrateur les 24 et 25/08/2026, ce qui est exactement l'événement qu'un journal
d'audit verrouillé doit rendre incontestable.

\textbf{Troisième seuil : cloisonner la donnée entre locataires.} C'est le seul des trois qui
demande plusieurs jours. La segmentation réseau qui occupait cette place a été appliquée du 11 au
25/08/2026 et mesurée par T20 ; ce qui reste relève du plan de la donnée : clé de locataire
obligatoire à l'ingestion, cloisonnement au stockage de l'entrepôt, instance de base dédiée au
locataire porteur de donnée irréparable --- conçue, chiffrée et ordonnancée sous D16
(\cref{subsec:isolation-donnees}), non appliquée. S'y ajoute le résidu réseau du correctif T14
(\cref{sec:non-conformite-t1}). Le dépôt pose lui-même ce seuil comme la condition d'accueil d'un
troisième client.

\section{Conditions de validité et réponse aux objectifs}
\label{sec:cloture}

\subsection{Conditions de validité des résultats}
\label{sec:limites}

Toute mesure porte les limites de son dispositif. Le \cref{tab:limites} recense celles qui
conditionnent l'interprétation des résultats, énoncées comme des bornes de validité et non comme
des excuses : chacune indique jusqu'où une conclusion peut être portée.

\begin{table}[!htbp]
\caption{Conditions de validité des résultats du chapitre}
\label{tab:limites}
\centering
\footnotesize
\setlength{\tabcolsep}{3pt}
\begin{tabular}{|P{3.9cm}|P{11.4cm}|}
\hline
\textbf{Condition} & \textbf{Portée qu'elle donne aux résultats} \\
\hline
Absence de trafic réel & Les résultats valent pour les scénarios exécutés, non pour une
exploitation réelle : ni taux de faux positifs ni seuil comportemental n'en est déductible \\ \hline
Attaques simulées, non adverses & Les tests valident les scénarios prévus, pas l'espace des
scénarios non prévus \\ \hline
Période d'observation courte & Les phénomènes lents --- dérive, saturation, coût --- restent hors
du champ d'observation ; c'est précisément l'objet du plan de validation continue
(\cref{sec:plan-validation}) \\ \hline
Isolation portée par les droits au plan du stockage & Les refus observés entre locataires
attestent le cloisonnement des identités, des secrets, des bases et --- depuis le refus croisé du
24/08/2026, T20 --- du réseau ; la conclusion ne s'étend pas au stockage, où l'instance de base de
données et l'entrepôt de supervision demeurent communs \\ \hline
Ressource d'infrastructure partagée & L'instance de base de données reste unique, seules les bases
et les rôles y étant séparés : point de défaillance unique que la campagne n'éprouve pas, et une
restauration restaure les deux bases ensemble. L'autre ressource commune, le connecteur d'accès
sans serveur, a été supprimée le 15/08/2026 \\ \hline
Compteurs incluant l'\textbf{auto-consultation} & Les volumes affichés par la vue d'ensemble
incluent le trafic généré par la consultation du tableau de bord lui-même : ils illustrent le
fonctionnement de la chaîne, ils ne mesurent pas une activité \\ \hline
Journal d'audit applicatif sans politique de purge & Aucune durée de conservation n'y est
déclarée : sa croissance est non bornée, et une ligne de politique suffirait, mais elle doit
exister avant tout engagement de service \\ \hline
Écriture non bloquante de la piste d'audit applicative & L'écriture est enveloppée dans un
traitement d'erreur silencieux : une panne ferait perdre la trace sans aucun signal, et les
conclusions tirées de cette piste valent sous réserve de sa complétude \\ \hline
\textbf{Dispositif de mesure mono-opérateur} & La reproductibilité est assurée par publication :
chaque protocole est écrit avec son critère d'acceptation, formulé avant exécution, de sorte qu'un
tiers puisse le rejouer et contredire le résultat (annexe~\ref{ann:protocoles-tests}) --- garantie
de méthode substituée à une revue croisée indisponible à cette taille \\ \hline
Annotation sémantique par une seule personne & L'évaluation comparative de l'enrichissement n'est
donc pas prononcée : le protocole prévoit une seconde annotation et un taux d'accord, et la grille
reste vide jusque-là (\cref{ann:evaluation-semantique}) \\
\hline
\end{tabular}
\end{table}

\subsection{Réponse aux objectifs du projet}

Les six objectifs mesurables du \cref{ch:cadre-existant} sont confrontés aux résultats obtenus
(\cref{tab:objectifs-ch6}). Une échelle binaire « atteint / non atteint » confondrait un objectif
dont le dispositif n'existe pas et un objectif dont le dispositif fonctionne sans être chiffré.
Trois niveaux cumulatifs sont donc employés --- \textbf{conçu} : traduit en exigences, en
architecture et en décisions datées ; \textbf{réalisé} : déployé et observé en fonctionnement ;
\textbf{mesuré} : une valeur relevée et datée établit le résultat.

{\footnotesize
\setlength{\tabcolsep}{3pt}
\begin{longtable}{|P{0.6cm}|P{2.6cm}|P{10.25cm}|P{1.6cm}|}
\caption{Réponse aux six objectifs mesurables du projet}
\label{tab:objectifs-ch6} \\
\hline
\textbf{\#} & \textbf{Objectif} & \textbf{Acquis, et portée de l'acquis} & \textbf{Niveau} \\
\hline
\endfirsthead
\hline
\textbf{\#} & \textbf{Objectif} & \textbf{Acquis, et portée de l'acquis} & \textbf{Niveau} \\
\hline
\endhead
O1 & Supprimer la confiance implicite du réseau & Contrôles d'identité et de périmètre déployés,
refus observés en recette et rejoués par test automatisé ; refus croisé entre locataires opposé et
journalisé le 24/08/2026 (T20). \emph{Portée} : frontière opposable à une charge compromise, non à
un opérateur capable de déployer ; deux charges conservent une sortie managée non gouvernée
(T14) & Réalisé, mesuré en partie \\ \hline
O2 & Rendre l'infrastructure reproductible & Socle décrit en code, construit puis reconstruit en
recette : la capacité est établie, la \emph{mesure} chronométrée relevant du plan de validation
continue, dont le versionnement du fichier de variables est le prérequis & Réalisé \\ \hline
O3 & Contrôler la chaîne de livraison & Portes bloquantes et \emph{éprouvées} : refus réel sur une
vulnérabilité critique (T8), fédération sans clé, aucune clé exportée (T7). \emph{Portée} : le
registre de confiance repose sur la construction, non sur un refus (T9) & Réalisé, mesuré en
partie \\ \hline
O4 & Détecter les événements de sécurité & Sept règles actives, chaîne complète démontrée le
19/08/2026, et publication de ce que le socle \emph{ne} détecte pas
(tableaux~\ref{tab:lacunes}, \ref{tab:limites} et \cref{sec:limites-filtrage}), comme l'objectif
l'exigeait. \emph{Portée} : couverture, faux positifs et délai médian non chiffrés ; la chaîne se
termine par un écran ; une interruption de la route de journaux ne produit aucune alerte (T12) &
Réalisé \\ \hline
O5 & Enrichir les alertes de façon déterministe & Candidats sémantiques produits, stockés et
consultables (\cref{fig:capture-K33}). \emph{Portée} : comparaison aux méthodes de référence non
prononcée et jointure au score d'incident laissée en écart assumé --- l'objectif admettait par
avance le résultat négatif ; débit plafonné à 200~détections par heure, au-delà desquelles la
fenêtre de deux heures écarte les plus anciennes sans erreur (T15, 22/08/2026) & Réalisé en
partie \\ \hline
O6 & Rester exploitable et soutenable & Restauration exécutée et chronométrée (32~min~45~s, perte
nulle, 314~lignes sur 314), retours arrière à 11,6~s et 16,5~s, haute disponibilité régionale
appliquée le 08/08/2026 en 11~min~2~s, plafond de charge mesuré le 22/08/2026, état de
l'infrastructure protégé par quatre refus le 20/08/2026. \emph{Portée} : bascule de zone non
provoquée, coût non ventilé & Réalisé, mesuré en partie \\
\hline
\end{longtable}
}

Aucun des six objectifs n'en reste au stade de la conception : les six dispositifs sont construits
et observés en fonctionnement, et ce qui n'est pas encore chiffré est de nature homogène --- durée,
volume, fréquence --- c'est-à-dire exactement le champ du plan de validation continue.

\section*{Conclusion du chapitre}
\addcontentsline{toc}{section}{Conclusion du chapitre}

\looseness=-1
Les vingt protocoles portent une preuve nommée, vérifiable et datée ; six portent en outre une
cadence de rejeu. Les deux plus fortes ne sont pas des captures mais un test automatisé qui tente
réellement l'écriture interdite sur les tables de preuves à chaque livraison, et un refus non
provoqué. Les non-conformités T1 et T14, conservées et expliquées, sont la contrepartie d'une
campagne dérivée de l'analyse de risque \emph{avant} la conception. Les points fragiles se nomment
aussi : couverture de détection non calculée, chaîne se terminant par un écran plutôt que par une
notification, deux charges dont la sortie réseau échappe au pare-feu, stockage commun aux deux
locataires. Restent les trois seuils --- alerter, verrouiller la rétention, cloisonner la donnée
--- qui séparent un socle qui fonctionne d'un service que l'on peut vendre, et la question ouverte
au \cref{ch:etat-art} : le gain de l'enrichissement sémantique justifiera-t-il son coût\,?

% Restauration des reglages de flottants par defaut
\setcounter{topnumber}{2}\setcounter{bottomnumber}{1}\setcounter{totalnumber}{3}
\renewcommand{\topfraction}{0.7}\renewcommand{\bottomfraction}{0.3}
\renewcommand{\textfraction}{0.2}\renewcommand{\floatpagefraction}{0.5}
\setlength{\LTpre}{\bigskipamount}\setlength{\LTpost}{\bigskipamount}
